\chapter{Riemann--Roch in Genus One}
\label{ch:riemann-roch}

For a divisor \(D\), define
\[
  L(D)=\{f\in\M(X): \divisor(f)+D\ge 0\}\cup\{0\}.
\]
Thus \(L(D)\) is the vector space of meromorphic functions whose poles are bounded
by \(D\).

\begin{theorem}[Riemann--Roch]
\label{thm:rr}
On a compact Riemann surface of genus \(g\),
\[
  \ell(D)-\ell(K-D)=\degdiv D+1-g,
\]
where \(K\) is a canonical divisor and \(\ell(D)=\dim_\C L(D)\).
\end{theorem}

We will use the theorem in its genus-one form.  On a complex torus, \(dz\) is a
nowhere-vanishing holomorphic \(1\)-form, so \(K\sim0\).  Since \(g=1\), the
formula becomes
\[
  \ell(D)-\ell(-D)=\degdiv D.
\]
If \(\deg D>0\), then \(L(-D)=0\), and therefore \(\ell(D)=\deg D\).

\begin{corollary}
\label{cor:ell-nO}
For \(T=\C/\Lambda\) and \(O=[0]\),
\[
  \ell(nO)=n\qquad(n\ge 1).
\]
\end{corollary}

This is the dimension count behind the fact that \(1,\wp,\wp'\) span \(L(3O)\).
The function \(1\) has no pole, \(\wp\) has a
double pole at \(O\), and \(\wp'\) has a triple pole at \(O\).  Since
\(\ell(3O)=3\), these three functions exhaust the space.

\begin{warning}
Riemann--Roch is a deep theorem in general.  In this book it is used as a
standard final-exam tool and interpreted carefully in genus one.  Claims that
require more algebraic geometry are explicitly labeled as such.
\end{warning}

\section*{Exercises}
\addcontentsline{toc}{section}{Exercises}

\exerciseintent{Use genus-one Riemann--Roch as a dimension calculator.}
\begin{exercise}
Using \(\ell(nO)=n\), compute \(\dim L(O)\), \(\dim L(2O)\), and
\(\dim L(3O)\) by substituting \(n=1,2,3\).
\end{exercise}

\exerciseintent{Show why \(L(O)\) contains no nonconstant function.}
\begin{exercise}
Split into two cases.  If \(f\in L(O)\) has no pole, use compactness to compute
that it is constant.  If \(f\) has a simple pole at \(O\), compute its total
pole order and compare with the elliptic-function pole-order obstruction from
Chapter 4.
\end{exercise}

\exerciseintent{Use pole orders to prove linear independence of \(1,\wp,\wp'\).}
\begin{exercise}
Suppose \(a+b\wp+c\wp'=0\).  Compare the pole orders \(0,2,3\) at \(O\) and use
the largest pole order term to compute successively that \(c=b=a=0\).
\end{exercise}

\exerciseintent{Combine membership, pole-order independence, and dimension to identify \(L(2O)\).}
\begin{exercise}
Use \(\ell(2O)=2\) and the pole order of \(\wp\) to prove
\[
  L(2O)=\Span_\C(1,\wp).
\]
Your proof should explicitly check membership of \(1,\wp\), linear independence
by pole order, and the dimension count.
\end{exercise}

\exerciseintent{Identify \(L(3O)\), the coordinate space behind the Weierstrass map.}
\begin{exercise}
Use \(\ell(3O)=3\) and the pole orders of \(\wp\) and \(\wp'\) to prove
\[
  L(3O)=\Span_\C(1,\wp,\wp').
\]
Your proof should explicitly check membership, pole-order independence, and the
dimension count.
\end{exercise}
